\itemhead{第3章 小结与复习 · 第1题（教材p.89）}{XJ-B1-C03-SU-P089-Y0854}
\label{XJ-B1-C03-SU-P089-Y0854}
{\footnotesize 湘教版教材p.89 / PDF第96页 / 第3章 小结与复习 / 第1题}
\field{原题}{初中已学习函数的概念，在本章，我们进一步用更加严谨的集合和对应语言来刻画函数，突出了函数概念的本质是一个数集到另一个数集的一种确定的对应关系，明确了构成函数的要素，并引入函数符号 $y=f(x)$。通过本章的学习，你对函数概念有什么新的认识和体会吗？}
\field{参考答案}{函数是非空数集之间“每个输入都有且只有一个输出”的对应；相同函数看定义域与对应关系，值域由它们共同决定；不要求存在统一解析式，也不要求不同输入有不同输出。}
{\small \field{关联知识}{F01 函数本质、构成要素与相同函数（教材帮p.77、78 / PDF86、87）}
\field{方法与易错点}{逐一核对每个允许输入是否有且仅有一个合乎目标集合的输出。 易错：把值域直接当作目标集合；只看解析式外观；认为常函数不是函数。}
\field{答案依据}{答案册未提供本栏目的独立对应答案；依据原题、教材概念补写参考回应。开放性总结不设唯一数值答案。}}
